\documentclass{article}
\usepackage{amsmath}
\usepackage{titling}

\usepackage[margin=0.5in]{geometry}

\title{Shoot on the Move Theory Proof\\
\large Charging Champions - 6560}

\author{Anish Kumar}

\begin{document}

\maketitle

\section{Theory}

By applying kinematics calculations based on odometry readings of the position and velocity of the robot and position of the target, vectors can be calculated and processed to result in a final target vector for velocity and direction of the shoot. This would then be translated through physical constants and some formulas to derive the RPM of the roller-compression system. Additional effects like air resistance and efficiency of spin transfer constants can be calculated post-implementation.

\section{Setup}

Goal: Get the RPM and turret angle to shoot at while moving at a velocity $\vec{v}$
\\
\\
There are 2 main vectors to keep track of (at the beginning):
\\
\\
let $\vec{r}$ be a "relative position vector", used to track the difference between the robot and target's positions regardless of vertical height 
\\
\\
let $\vec{v}$ be a yagsl-determined vector to show the robots velocity in 2 dimensions, also disregarding vertical height
\\
\\
These will not be combined directly, as they are position and velocity, and instead operations will be performed to translate the position vector into a velocity vector.
\\

\section{Kinematics}

Imagine a triangle in a 2d plane, similar to a standard kinematics projectile motion problem with vectors for velocity and it's components forming a right triangle with an angle $\theta$
\\
Let the vectors on this triangle be made of 
\[
\vec{l} , \vec{l}_x , \vec{l}_y
\]
\\
By solving like a normal projectile motion question, we can separate the vector into componenets and solve from there
\[y(t) = y_0 + \|\vec{l}\|sin(\theta) t - \frac{1}{2}gt^{2}\]
\\
And

\[x(t) = \|\vec{l}\|cos(\theta) t\]
\\
But, drag aplies to both the $x$ and $y$ functions
So, using a standard calculation for drag force:

\[F_{drag} = \frac{1}{2} \cdot \rho \cdot A \cdot C_d \cdot v^2\]
\\
Density of air is $\approx{1.225}$\\
Contact area = $\frac{\pi d^2}{4} = \pi \cdot 0.005625$m $\approx 0.01767$


%recompute all associated values or approximations because of the pi{d/2}^2 thing 


\[F_{drag} = \frac{1}{2} \cdot 1.225 \cdot 0.01767 \cdot C_d \cdot \vec{l}^2 \approx 0.2886345 \cdot C_d \cdot \|\vec{l}\|^2\]
\\
Substituting into $x(t)$
\[x(t) = \|\vec{l}\|cos(\theta) t - \frac{1}{2m}(F_{drag})t^2\]
\[x(t) = \|\vec{l}\|cos(\theta) t - \frac{1}{2m}(0.2886345 \cdot C_d \cdot \|\vec{l}\|^2)t^2\]
\[2mx = mt(2\|\vec{l}\|cos(\theta) - (0.2886345 \cdot C_d \cdot \|\vec{l}\|^2)t)\]
\[-(0.2886345C_d\|\vec{l}\|^2)t^2+2mt\|\vec{l}\|cos(\theta)-2mx = 0\]
\\(Holy what is this)
\[\|\vec{l}\| = \frac{-2mtcos(\theta) \pm \sqrt{(-2mtcos(\theta))^2 - 4(-0.2886345C_dt^2)(-2mx)}}{2(-0.2886345C_dt^2)}  \]
\\
\\
Using an angle 45\textdegree \,as the value for theta (based on alphabot for now)
\[\|\vec{l}\| = \frac{-mt\sqrt{2} \pm \sqrt{(-mt\sqrt{2})^2 - 4(-0.2886345C_dt^2)(-2mx)}}{2(-0.2886345C_dt^2)}\]
\[ = \frac{-mt\sqrt{2} \pm \sqrt{2m^2t^2 - 2.309076C_dm^2t^2x}}{-0.577269C_dt^2}\]
\[ = mt \cdot \frac{-\sqrt{2} \pm t(\sqrt{2 - 2.309076C_dx})}{-0.577269C_d}\]
\\\\
Take only the $+$ solution because we want the highest arc:
\[ = mt \cdot (\frac{2.44983458729}{C_d} + \frac{t(\sqrt{2 - 2.309076C_dx})}{-0.577269C_d})\]
\\
Drag Coefficients can be determined experimentally, but seeing as we don't have an easy way to experimentally do so, we can estimate it by using Reynold's numbers. 
\\\\
$C_d$ is calculated as a piecewise function of the object's Reynold's number ($Re$), 

\[Re_{fuel} = \frac{\rho\|\vec{l}\|D}{\mu}\]
\\\\
$\rho$ is the density of air, $D$ is the diameter of the sphere, $\mu$ is the viscosity of air.

\[Re_{fuel} = \frac{1.225 \cdot \|\vec{l}\| \cdot 0.15}{1.8 \cdot 10^{-5}}\]

\[\approx 10,208.33\|\vec{l}\|\]
\\
And with exit velocity being in the range from 5 to 20 mph,

\[Re_{fuel} \approx 10,208.33(2.24 \,\, to \,\, 9.94 \frac{m}{s}) \approx 2 \cdot 10^4 \,\, to \,\, 9 \cdot 10^4\]
\\
Since this value is in the subcritical drag range, below drag crisis, $C_d$ can be approximated as $0.47$

\[\|\vec{l}\| = mt \cdot (\frac{2.44983458729}{C_d} + \frac{t(\sqrt{2 - 2.309076C_dx})}{-0.577269C_d}) \approx mt \cdot (5.212414 + \frac{t(\sqrt{2-1.08526572x})}{-0.27131643})\]
\[\approx 5.212414mt +\frac{mt^2(\sqrt{2-1.08526572x})}{-0.27131643}\]
\\\\
Fuel mass according to the game manual is $m_{fuel} \approx (0.203 \,\, to \,\, 0.227kg)$ so can be used as $\approx 0.215 \,kg$
\\\\
\[\boxed{\|\vec{l}(x,t)\| \approx 1.12t - 0.8t^2\sqrt{2-1.085x}}\]
\\\\
Substituting back into $y(t)$

\[ y = \frac{\|\vec{l}\|t}{\sqrt{2}} - \frac{(g-\frac{F_{drag}}{m})t^2}{2}\]
\[2 \sqrt{2}y = 2t\|\vec{l}\| - \sqrt{2}t^2\cdot(g-\frac{0.2886345 \cdot C_d \cdot \|\vec{l}\|^2}{m}) \approx 2t\|\vec{l}\| -\sqrt{2}t^2 \cdot (g-0.630968\|\vec{l}\|^2)\]
\\
Solve for $\|\vec{l}(y,t)\|$

\[-\sqrt{2}t^2 \cdot (g-0.630968\|\vec{l}\|^2) + 2t\|\vec{l}\| - 2 \sqrt{2}y = 0\]
\[\|\vec{l}\| = -\frac{1.122}{t} + \frac{\sqrt{4 + 10.08y + 5.04g}}{1.782t}\]
\[\boxed{\|\vec{l}(y,t)\| = -\frac{2 + \sqrt{53.4424 + 10.08y}}{1.782t}}\]
\\
Then solve for time of flight:

\[-2 - \sqrt{53.4424 + 10.08y} = 1.782t(1.12t - 0.8t^2\sqrt{2-1.085x})\]
\\
The nature of this function means it is irreducible, even through Newton-Raphson method. However, time of flight can be implicitly defined as a function of $x$ or $y$,
\[\boxed{t(x,y) = -1.4t^3\sqrt{2-1.1x} + 2t^2 + \sqrt{53.44 + 10.1y} + 2 = 0}\]
\\
When actually running, $x$ and $y$ would be inputted into $t(x,y)$. This would result in a cubic, irreducible function of solely $t$, meaning Newton-Raphson method can be used to approximate the roots of the function.
From there, the value $t$ and the $x$ could be substituted into $\|\vec{l}(x,t)\|$ to finally determine the exit velocity. 

\newpage

\section{Vector Processes}

Now that there is a vector for magnitude in that plane, we need to work in the ground plane for direction calculations:
\\
\\
First, convert $\vec{l}$'s magnitude to a ground-relative vector, aka $\vec{l_x}$, and call it $\vec{g}$ to store ground velocity of the ball.

\[\|\vec{g}\| = \|\vec{l}\|cos(\theta) = \frac{x}{t} = x\sqrt{\frac{g}{2(x-y)}}\]
\\\\
Since the stationary ground velocity of the ball would retain the heading of the position vector $\vec{r}$,

\[arg(\vec{g}) = \phi = atan2(\vec{r_y},\vec{r_x})\]
This would mean that 0\textdegree is upfield and $\phi$ would increase counter-clockwise from that
\\\\
Then break it into components:

\[\vec{g} = \|\vec{l}(x,t)\|cos(\phi)\hat{\imath} + \|\vec{l}(y,t)\|sin(\phi)\hat{\jmath}\]
\\
This vector $\vec{g}$ is the field relative target velocity vector. 
\\
Next we have to determine the robot's field-relative velocity $\vec{v}$
Given a gyro reading of $\beta$, $\vec{v}$ can be calculated

\[\vec{v} = (\vec{v_x}cos(\beta) - \vec{v_y}sin(\beta-\frac{\pi}{2}))\hat{\imath}\]

\end{document}